\section{Coarse entry from contact}\label{v9:entry:section}
Fix an attained pair with $a\ge0$, $n\ge N=42$, and true predecessor
$C_{n-1}\le c=\CE$. All bounds below concern this one marginal and raw
threshold. They pass through two quadratic and four signed updates;
their finite coverage conditions remain explicit.

\subsection{Support before higher moments}
\begin{theorem}[Initial entry]\label{v9:entry:initial}\label{v9:entry:initial-support}
Assume Shevtsova's external structural inequality \cite{Shevtsova2013}
\begin{equation}\label{v9:entry:structural}
 \sqrt n\Delta_n(P)\le\frac{3031}{10000}(\rho+323/500).
\end{equation}
Then $\rho<7361/4000$, $|\beta|<321/500$ and
$\supp P\subset(-96/25,96/25)$, uniformly in $H$.
\end{theorem}
\begin{proof}
Since $c\rho\le\sqrt n\Delta_n(P)$ and $c>819/2000$, division gives
\[
 \rho<\frac{(3031/10000)(323/500)}{819/2000-3031/10000}=\frac{7361}{4000}.
\]
Signwise Cauchy gives $\rho\ge(1+\beta^2)/\sqrt{1-\beta^2}$,
increasing in $|\beta|$. At $r=1841/1000$, $b=321/500$,
\[
 (1+b^2)^2-r^2(1-b^2)=\frac{93288849}{50000000000}>0.
\]

Put $v=n-1$, $s=\sqrt{n/v}$, $F_v(x)=\Phi(x/\sqrt v)$ and
$e_n=n(s-1)-1/2$. Contact and the predecessor imply
\begin{equation}\label{v9:entry:contact}
 0\le H_n(H,t)+c\{-|t|^3+\tfrac32\rho t^2+3\beta t+\rho e_n\}-na
 \quad(t\in\supp P),
\end{equation}
where
\[
 H_n/n^{3/2}=F_v(H-t)-F_{v+1}(H)+tF'_{v+1}(H)
                         -(t^2-1)F''_{v+1}(H)/2.
\]
The extra $-na$ is nonpositive; no unknown-score comparison is used.
Rationalization gives
$ne_n=s^2(2s+1)/(2(s+1)^2)<793881/2073680<48/125$ and
$s<S:=81/80$, hence $e_n<\eta:=8/875$; the rationalized expression
increases for $s\ge1$.

The heat identity $\partial_vF_v=F_v''/2$, spatial Taylor expansion to
order four and variance expansion to order two give, for $u\ge1$,
\begin{equation}\label{v9:entry:heat}
 H_n(H,\pm u)/s^3
 =\pm A(1-z_v^2)\phi(z_v)+B\phi'''(z_v)+\mathcal E,
\end{equation}
where $z_v=H/\sqrt v$, $A=(u^3-3u)/6$, $H_4=u^4-6u^2+3$,
$B=H_4/(24\sqrt v)$ and
\[
 |\mathcal E|<E(u):=u^5/4100+3u/820+17(3u^2-2)/78720.
\]
The four remainders after division by $s^3$ are bounded by
\[
 \frac{u^5\|\phi^{(4)}\|_\infty}{120v},\quad
 \frac{\|\phi^{(5)}\|_\infty}{48v^{3/2}},\quad
 \frac{u\|\phi^{(4)}\|_\infty}{8v},\quad
 \frac{|u^2-1|\|\phi^{(5)}\|_\infty}{16v^{3/2}}.
\]
Gaussian inversion bounds these derivative norms by $6/5$ and $51/20$;
use $v\ge41$, $\sqrt v>6$. No higher moment of $P$ enters.
For $A,B>0$, $B\le A/4$, the leading expression is at most
$\phi(0)(A+9B^2/(4A))$ for both signs: for the positive sign on
$|z_v|\le1$, complete the square in $A-Az_v^2+3B|z_v|$; outside it the
first term is nonpositive. For the negative sign use
$2e^{-3/2}A<A/2$ and $\|\phi'''\|_\infty<2\phi(0)$.
The latter follows from $3x-x^3\le2$ on $[0,\sqrt3]$ and
$x^3e^{-x^2/2}\le3\sqrt3e^{-3/2}<2$ beyond it.

Freeze the preceding $r,b$, and set
\[
 Q(u)=-u^3+3ru^2/2+3bu+\eta r,
\]
\[
 P(u)=\tfrac25S^3\{A^2+H_4^2/(256\cdot41)\}
                      +S^3E(u)A+(819/2000)Q(u)A.
\]
At $R=96/25$ we have $A>0$, $0<H_4<36A$, $Q(R)<0$ and
\[
 -P(R)=\frac{2214258283536204102723}{10250000000000000000000}>\frac15.
\]
Contact is therefore impossible at either boundary. For the whole rays,
\[
 |\partial_tH_n(H,t)|<\tfrac25(t^2/2+1/2+e_n).
\]
Indeed, with $y=(H-t)/\sqrt n$, $z=H/\sqrt n$, the derivative is
$n\phi(z)+Ht\phi(z)-ns\phi(sy)$. The spatial second-order error costs
$\phi(0)t^2/2$, while $\partial_s[s\phi(sy)]=-\phi''(sy)$ costs
$\phi(0)n(s-1)$. Put
$L(u)=\tfrac25(u^2/2+1/2+\eta)+(819/2000)Q'(u)$.
At $R$, all $Q,Q',Q'',L,L',L''$ are negative. As $Q'''=-6$ and
$L''=2/5-6(819/2000)<0$, the signs persist outwards, so both actual
radial upper expressions decrease from negative boundary values.
\end{proof}
The same heat/ray argument licenses a later radius in
$\{2.50,2.51,\ldots,3.84\}$: freeze the cell bounds $r,b$, then check
$A>0$, $0<H_4<36A$, $Q(R),P(R)<0$ and the six ray signs before
using that radius in any higher-moment bound.

\subsection{The common Fourier input}\label{v9:entry:moments}\label{v9:entry:fourier}
After acquiring support, put $M_j=\E|X|^j$. The measure
$d\mu_0=(x-1)^2(x+1/2)dP$, $x=|X|$, has mass $\rho-1$, and
\[
 M_4-1=\int\frac{(x+1)^2}{x+1/2}d\mu_0,\qquad
 M_5-1=\int\frac{x^3+2x^2+3x+3/2}{x+1/2}d\mu_0.
\]
The ratios increase on $x\ge0$; write $k_4,k_5$ for their values at $R$.
Consequently
\begin{equation}\label{v9:entry:mom}
 M_4\le\min\{R\rho,1+k_4(\rho-1)\},\quad
 M_5\le\min\{RM_4,1+k_5(\rho-1)\},\quad M_6\le RM_5.
\end{equation}
Also $\gamma^2\le M_4-1$ by $\E(X^2-1-\gamma X)^2\ge0$.
For two iid copies, $\mathcal E=\{(X-X')^2\ge9/25\}$ satisfies
\begin{equation}\label{v9:entry:event}
 \Pp(\mathcal E)\ge\alpha:=\frac{(2-9/25)^2}{2M+6},\quad
 |X-X'|\ge3/5\text{ on }\mathcal E,\quad
 \E[(X-X')^2\mid\mathcal E]\le2/\alpha,
\end{equation}
where $M\ge M_4$. Cauchy proves this from
$\E[(X-X')^2;\mathcal E]\ge2-9/25$ and
$\E(X-X')^4=2M_4+6$; the last bound uses $\E(X-X')^2=2$.

Appendix~\ref{v9:fourier:section} proves the complete compensated
Fourier bound, including atoms and the comparator tail. Its input is
an acquired radius $R$, ceilings $M\le11.1$, $g\le1.85$ for $M_4,|\gamma|$,
a sixth-moment ceiling, and
$|(3-M_4)/24|\le A_4$, $|(10\gamma-\E X^5)/120|\le Z_5$.
Use the complete LOW triple $L_0,L_H,L_2$ of
Proposition~\ref{v9:fourier:low}, with the sharper fifth-order coefficient
proved there. An actual event of mass at least $\alpha$, separation at
least $d_E$ and conditional second moment at most $M_E$ supplies the
six-term HIGH bound $K_H$ of Proposition~\ref{v9:fourier:high}, under all
its damping and all-index gates. Insert each event option into all six
terms before minimization. The output is
\begin{equation}\label{v9:entry:fourier-output}
 |E_n|\le K_F+s_F|H|,\qquad
 K_F=\tfrac{16}{25}(L_0+K_H+L_2/\sqrt N),\quad s_F=\tfrac{16}{25}L_H.
\end{equation}
It is uniform in actual $n\ge N$, even when the law and threshold vary.

\subsection{Coarse coefficients and two passes}\label{v9:entry:coarse-coefficients}
The constant coefficient of $c$ in \eqref{v9:core:R} is
\[
 A^{\rm law}=3\beta(J+\gamma)/2+3(\rho J_1-J_2)
                         +\rho M_4/2-M_5/4-\rho/2.
\]
Without a central gap one still has
\begin{equation}\label{v9:entry:gap-free}
 -M_4/6\le J_1\le M_4/2,\qquad-3J_2-M_5/4\le0,\qquad J\le\rho+1.
\end{equation}
For a gap $(u,v)$ of length $h$, the first pair follows from
$\int(t-u)t+\frac12\int t^2=hv^2/2$ and
$\frac32\int t^2-\int(t-u)t=h(2u+v)^2/6$; sum against its weight
and use $\int t^2\tau=M_4/3$. On negative gaps reflect and use
$\frac13\int_u^vx^3dx-\int_u^v(v-x)x^2dx=u^3(v-u)/3\ge0$ for the
second bound; positive pieces improve it. The last uses $J\le\rho+2m$.
For $1\le r\le\rho$, $R\le6$, choose
\begin{equation}\label{v9:entry:Gamma}
 \Gamma(R,r)\ge\max_{r\le x\le R}
 \left\{-x^2/4+(R+r)x/3-Rr/2+r^3(2R-r)/(12x^2)\right\}_+.
\end{equation}
Then
\begin{equation}\label{v9:entry:coarse-gap-charge}
 3(\rho J_1-J_2)\le3\Gamma(R,r)(\rho-\gamma)/2+3\rho^2J/8.
\end{equation}
On a negative gap discard $|t|<\rho$, whose charge is nonpositive.
Extend the rest to $\int_r^v(R-x)x(x-r)dx$; dividing by $v^2$ gives
\eqref{v9:entry:Gamma}. Tonelli charges its bound against
$\E X_-^3=(\rho-\gamma)/2$. On positive pieces use
$t(\rho-t)\le\rho^2/4$ and $\sum\tau\int(t-u)=J/2$.
To enclose $\Gamma$, split $[r,R]$ into 64 closed intervals, maximize the
quadratic at its clipped vertex $2(R+r)/3$ and the positive reciprocal
term at the left endpoint. Alternatively $\Gamma\le5/2$: at $r=1,R=6$
this is $3x^4-28x^3+66x^2-11>0$ on $[1,6]$, whose endpoint values are
$30,205$ and interior critical values $(85\mp35\sqrt5)/2>0$.
The integral increases with $R$ and decreases with $r$.

On $\rho\in[l,u]$, $|\beta|\le b$, $M_4\le M$, $|\gamma|\le g$ and
$0\le\rho-\gamma\le\Delta$, the square
$V=M_5-2\rho M_4+\rho^3\ge0$ gives two complete constant bounds:
\begin{align}\label{v9:entry:coarse-A}
 A_1(\Delta)&=3b(2u+1)/2+3\Gamma(R,l)\Delta/2
                       +3u^2(u+1)/8+u^3/4-l/2,\notag\\
 A_2&=3b(2u+1)/2+2uM-l/2.
\end{align}
For $K_1=\E[X|X|^3]-3\beta-3\rho\gamma/2$, intersect the applicable bounds
\begin{gather*}
 K_{\rm rad}=3ub+\frac{R(R+2-3l/2)}{R+1/2}(u-1),\quad
 K_{\rm univ}=6u-3u^2/2+3b,\\
 K_R=Ru/2+\max_{r=l,u}|R/2-3r/2|g+3b.
\end{gather*}
The first follows from
$K_1+3\rho\beta=\int\operatorname{sgn}(X)x(x+2-3\rho/2)/(x+1/2)d\mu_0$.
This ratio is convex in $x$, decreases in $\rho$, and its negative part
is at most $31/40$ for $\rho\le1.85$; require the endpoint ceiling to
exceed that number. The second uses
$|\E[X|X|^3]-3\gamma|\le3\rho$ for $R\le6$, $\rho\le2$; the third
uses $0\le\E X_\pm^4\le R(\rho\pm\gamma)/2$. Finally
$|\varphi\Omega|\le v_H:=\frac25\max\{1/2,2M/3-1/2\}$ for
$\Omega=J_1+(M_4-3)/6$.

\begin{lemma}[Signed fourth-moment charge]\label{v9:entry:signed-K}
Put $h_*=1+b$, $a(r)=(3r-R)/2$. Another bound for $|K_1|$ is
$\max_{r\in[l,u]}\{Rr-3r^2/2+3b+(1+b)a(r)^2/2\}$.
If
\[
 h_*\ge a(u),\quad3+2a(l)h_*-3h_*^2/2\ge0,\quad R-3u+3h_*^2/2\ge0,
\]
intersect it with $Ru-3u^2/2-3+3h_*+a(u)h_*^2-h_*^3/2$.
\end{lemma}
\begin{proof}
For $\epsilon=\pm1$, put $h=1-\epsilon\beta>0$, $\delta=\rho-\epsilon\gamma$.
Signwise Cauchy and $m\le1/2$ give $\delta\ge h^2$; fourth moments give
\[
 \epsilon K_1\le R\rho-3\rho^2/2-3+3h+a(\rho)\delta-\delta^2/(2h).
\]
Unrestricted square completion proves the first bound. Under
$\delta\ge h^2$, the last three terms maximize to $h(3+a^2/2)$ for
$h\le a$, and $3h+ah^2-h^3/2$ for $h\ge a$. They agree at $a$.
The first is increasing; the second has derivative $3+2ah-3h^2/2$,
nonnegative on $[0,h_*]$ by the second gate. The first gate selects its
branch and the third makes its remaining $\rho$ derivative nonnegative.
Evaluate the first quadratic at endpoints and its eligible concave vertex.
\end{proof}
Call the applicable minimum $K$; no sign of $\beta\gamma$ is assumed.

Use 24 closed cells $[r_i,r_{i+1}]$, $r_i=1+3361i/96000$, and
\begin{gather*}
 c_-=.4097321837023<c<c_+=.4097321837025,\\
 \phi_l=997/2500<\phi(0),\qquad\sigma_l=3.162277660<\sqrt{10}.
\end{gather*}
The alternating Machin bounds and rational squares of
Lemma~\ref{ms:gaussian-constants} license these constants. Put $\sigma_u=3.162277661>\sqrt{10}$.
In pass 1 take $A=\min(A_1(2u),A_2)$,
$f(t)=tK+v_H+s_F$, $q(t)=tl-g/5>0$.
With overlines denoting upward rounding to $10^{-10}$, set
\begin{align}\label{v9:entry:first-pass}
 T_1&=\overline{\max_{t=c_-,c_+}\{K_F+tA+f(t)^2/(4q(t))\}},\notag\\
 \eta_1&=(c_-l-T_1/N)/2,\quad
 U_1=\min\{T_1,\overline{\max_{t=c_-,c_+}
                   [K_F+tA+f(t)^2/(4(q(t)+\eta_1))]}\},\notag\\
 d_1&=\overline{6U_1/(N\phi_l)}.
\end{align}
Check $0\le T_1<Nc_-l$ before using $\eta_1$. The raw quadratic is affine
in $c$ at each $|H|$, so maximize at endpoints before completing squares.
Corollary~\ref{v9:core:phase-price} then gives $T_1,U_1,Nd_1$.
For pass 2 the moment circle gives
$D\ge\sigma\rho-\gamma-3\sqrt{\rho^2-\gamma^2}$. Put $a_0=\sigma_l-d_1/l$ and
\begin{equation}\label{v9:entry:circle-inversion}
 y_l\le Y(a_0):=\begin{cases}-1,&a_0\le-1,\\
 (a_0-3\sqrt{10-a_0^2})/10,&-1<a_0\le\sqrt{10}.
 \end{cases}
\end{equation}
This smaller root inverts the increasing branch of
$y+3\sqrt{1-y^2}$ on $[-1,1/\sqrt{10}]$; $a_0-y\ge0$ prevents the
false squared branch. Round the root downward. Thus
\begin{gather*}
 \gamma\ge g_l:=\max\{-g,\sigma_l l-3-d_1,\min(ly_l,uy_l)\},\\
 \Delta=\min\{2u,(1-\sigma_l)l+3+d_1,u-g_l\},\quad
 \phi_-=\phi_l(1-T_1/(Nc_-l)),\\
 A=\min\{A_1(2u),A_1(\Delta),A_2\},\quad
 q(t)=\max\{tl-g/5,tl+\tfrac12\min(\phi_-g_l,(2/5)g_l)\}.
\end{gather*}
Here $2m\le1$ gives the additional bounds and both alternatives for $q$
have slope $l$. Repeat the raw step, intersect with $T_1$, price phase
only then, and intersect the resource with $U_1$ and $d_2$ with $d_1$.
\begin{condition}[Two coarse passes]\label{v9:entry:two-pass-condition}
All coefficient, support, Fourier, denominator and strict raw gates above
hold on all 24 cells, with outward ceilings
\[
\begin{array}{c|rrr|c}
 &\max T&\max U&\max d&\text{strict gap gates}\\\hline
1&13.2818901598&10.9330610254&3.9164138937&23\\
2&11.5164967977&9.9790315907&3.5746638454&24
\end{array}
\]
and $d_2<(\sigma_l-1)l$ in every second-pass cell, with minimum margin
\[
 \frac{526024602543}{1600000000000}>0.
\]
\end{condition}
This is a finite exact obligation. Zero in the support would give
$J\le\rho$ and $D\ge(\sigma-1)\rho$, contradicting its gap test.
Centering therefore acquires an actual central gap $(-A,B)$, $A,B>0$.

\subsection{Signed refinements of one physical law}\label{v9:entry:signed-moments}
Put $a_\pm=(1\pm\beta)/2$, $b_\pm=(\rho\pm\gamma)/2$. Signwise Cauchy gives
\begin{gather*}
 m^2\le a_+a_-,\quad mb_\pm\ge a_\pm^2,\quad
 m\ge\max\{(1+\beta^2)/(2\rho),a_+^2/b_+,a_-^2/b_-\},\\
 \frac{(1+\beta)^{3/2}}{\sqrt{1-\beta}}-\rho\le\gamma
 \le\rho-\frac{(1-\beta)^{3/2}}{\sqrt{1+\beta}}.
\end{gather*}
Both signs have positive mass and first moment, hence $|\beta|<1$.
Projection of $X|X|$ onto $(X,|X|)$, whose Gram matrix is
$\bigl(\begin{smallmatrix}1&\beta\\\beta&1\end{smallmatrix}\bigr)$, gives
$M_4\ge\rho^2+(\gamma-\rho\beta)^2/(1-\beta^2)$.
Integrating $(1,|X|)^{\mathsf T}(1,|X|)$ against
$(R-|X|)dP$, $|X|(R-|X|)dP$ and $|X|^2(R-|X|)dP$ on each sign gives
positive-semidefinite matrices, whose determinants imply
\begin{gather*}
 2m+\sum_\pm\frac{(Rm-a_\pm)^2}{D_\pm}\le R,
 \quad M_4\le R\rho-\sum_\pm\frac{D_\pm^2}{Rm-a_\pm},
 \quad D_\pm=Ra_\pm-b_\pm,\\
 \E X_\pm^5\le R\E X_\pm^4
                   -\frac{(Rb_\pm-\E X_\pm^4)^2}{D_\pm}.
\end{gather*}
Use corrections only at positive denominators. A positive numerator
with zero denominator excludes only by its original determinant
inequality; zero-over-zero corrections are dropped. The same $m$ is
retained on both signs. Integrating $(x-a)^2(R-x)\ge0$ gives also
$2m\le(R-\rho+2a+Ra^2)/(a(2R+a))$; use $a=2/10,\ldots,12/10$.

For $\delta=\gamma-\rho\beta$, three independent losses improve
$\rho M_4/2-M_5/4\le\rho^3/4$:
\[
 P_1^*=\frac{(\rho+\gamma)\delta_+^2}{8(1+\beta)^2}
       +\frac{(\rho-\gamma)(-\delta)_+^2}{8(1-\beta)^2},\quad
 P_2^*=\frac{\delta^4}{4\rho(1-\beta^2)^2},\quad P_3^*=\frac{\delta^2}{8m}.
\]
For the first, signwise Cauchy gives $c_4\ge b_3^2/a_2$, $f_5\ge c_4^2/b_3$,
so $\rho c_4/2-f_5/4\le\rho^2b_3/4-(c_4-\rho b_3)^2/(4b_3)$;
minimize the square over that half-line. The second follows from
$M_4^2\le\rho M_5$ and the projection bound. For the third,
$V=\E[|X|(|X|^2-\rho|X|)^2]$ and weighted Cauchy give
$\delta^2\le2mV$. Subtract their maximum, never their sum.

The physical decomposition~\eqref{ms:defect-decomposition} gives
$T_++T_-\le D/(\sigma-1)$ for $T_+=b_+-mB^2$, $T_-=b_--mA^2$.
Also $J\le\rho+2mAB$ gives $6mAB\ge(\sigma-1)\rho-D$.
On $\rho\ge l$, $D\le d<(\sigma_l-1)l$, $m\le m_u$,
\begin{equation}\label{v9:entry:gap-distance}
 (A+B)^2\ge\{l-d/(\sigma_l-1)+((\sigma_l-1)l-d)/3\}/m_u>0.
\end{equation}
Retain $A\le a_-/m$, $B\le a_+/m$ and $D\ge\sigma\rho-\gamma-6m$.
In particular
\begin{equation}\label{v9:entry:mean-contraction}
 m\ge m_*:=\max\{0,(\sigma_l l-\gamma_+-d)/6\},\qquad
 |\beta|\le\sqrt{1-4m_*^2}.
\end{equation}
An $m_*>1/2$ excludes the cell. Intersect beta afresh in each gamma
slice before recomputing dependent moments. The stronger outer-gap
inequality gives
$D\ge(\sigma-1)b_++6a_+B-(\sigma+5)mB^2$.
If $D<(\sigma-1)b_+$, its concave quadratic forces $B$ above the positive
root; the analogue uses $A,a_-,b_-,\sigma+1,\sigma+7$.
Round positive coefficients down, negative coefficients up and roots down.

The actual event $XX'<0$ has probability $2pq$ and separation $A+B$.
Cauchy gives
$p\in[2m_-^2/(1+\beta_+),1-2m_-^2/(1-\beta_-)]$ and
$2pq\ge(1+\beta^2)^2/(2\rho^2)$. Exactly,
\[
 \E[(X-X')^2\mid XX'<0]
 =\frac{1+4m^2+\beta}{2p}+\frac{1+4m^2-\beta}{2q}.
\]
This increases with $m$, is affine in $\beta$ and convex in $p$, with
positive numerators. Use its corner maximum or $2/\alpha$, whichever is
smaller, with
$\alpha=\max\{(1+\min|\beta|^2)^2/(2u^2),2\min_{p=p_-,p_+}pq\}$.
The downward root of \eqref{v9:entry:gap-distance} then licenses another
complete HIGH option, minimized against the universal event.

For the signed integral charge put
\[
 P_1=B^3/3+AB^2/2-A^3/6,\quad P_2=B^4/4+AB^3/3-A^4/12,\quad G=\rho P_1-P_2.
\]
The central terms are $J_{1,c}=mP_1$, $J_{2,c}=mP_2$.
Since $\partial_BG=B(A+B)(\rho-B)$,
\[
 G\le G(A,\rho;\rho)=\rho^4/12+\rho^3A/6-\rho A^3/6+A^4/12.
\]
It increases with $A$ (derivative $(\rho-A)^2(\rho+2A)/6$) and is convex
in $\rho$, so use $A_u$ and both rho endpoints. Positive outer gaps
satisfy $(t-a)(\rho-t)\le(\rho-B)_+^2/4$; negative outer gaps cost at
most $(R-\rho)^2/4$ after extracting $|t|$. Integration gives
\[
 3(\rho J_1-J_2)_o\le3\{(\rho-B)_+^2T_++(R-\rho)^2T_-\}/8.
\]
The joint budget is $D-mQ$, where
$Q=(\sqrt{\sigma+1}A-\sqrt{\sigma-1}B)^2$; intersect it with
$T_+\le b_{+,u}-m_lB_l^2$, $T_-\le b_{-,u}-m_lA_l^2$.
For $\Omega$ use $-RT_-/2\le J_{1,o}\le RT_+/2$.
The minimum of $P_1$ is at $B_l$ and an $A$ endpoint; the maximum is
at $B_u$ and $A$ clipped to $B_u$. Multiply a nonnegative polynomial bound
by $m_u$ for an upper bound and $m_l$ for a lower bound, reversing the
choices for a negative bound. Fourth-semimoment intervals similarly
enclose $K_1$. Retain the full signed pair $(K_1,\Omega)$.

\subsection{Ordered update and closed coverage}\label{v9:entry:ordered-update}
Each cell carries its radius, moments, raw ceiling $T$ and weighted
invariant $nD\le Nd$. Intersect all preceding necessary conditions
outwardly; monotone bisection for $m$ retains its upper endpoint. Take
\[
 J_- =\max\{0,((3+\sigma_l)l-\gamma_+-d)/3\},\quad
 J_+=\min\{u+2m_+,((3+\sigma_u)u-\gamma_-)/3\}.
\]
Let $B_*$ be the corner maximum of $3\beta(J+\gamma)/2$, $G_*$ the smaller
coarse/central-plus-outer charge and $P_*$ a lower bound for
$\max(P_1^*,P_2^*,P_3^*)$. Then
\begin{equation}\label{v9:entry:signed-constant}
 A=\min\{B_*+G_*+u^3/4-l/2-P_*,B_*+2uM-l/2\}
\end{equation}
encloses the complete constant. Let $[k_-,k_+]$, $[\omega_-,\omega_+]$
enclose $K_1,\Omega$. Under $0\le T<Nc_-l$, put
$\phi_-=\phi_l(1-T/(Nc_-l))>0$ and
\begin{align}\label{v9:entry:signed-quadratic}
 B_t&=K_F+tA,\quad
 F_t=\max_{k=k_\pm,\,\omega=\omega_\pm,\,\psi=\phi_-,2/5}
                        |tk+\psi\omega|+s_F,\notag\\
 Q_t&=tl+\tfrac12\min\{\phi_-\gamma_-,(2/5)\gamma_-\}.
\end{align}
Thus $\mathcal R_c\le B_c+F_c|H|-Q_cH^2$. Successive affine maximization
licenses the corners; convexity in $c$ then licenses $c_-,c_+$ even when
$A<0$. Negative $\gamma_-$ uses the upper density endpoint.
\begin{lemma}[Ordered step]\label{v9:entry:signed-step}
If $Q_t>0$ at both endpoints, with all ceilings rounded upwards, set
\begin{gather*}
 T'=\min\{T,\max_{t=c_-,c_+}(B_t+F_t^2/(4Q_t))\},\quad
 \eta=(c_-l-T'/N)/2,\\
 U'=\max_{t=c_-,c_+}(B_t+F_t^2/[4(Q_t+\eta)]),\quad
 d'=\min\{d,6U'/(N\phi_l)\}.
\end{gather*}
Negative $T'$ or $U'$ excludes the cell. Otherwise
$\mathcal R_c\le T'$, $\mathcal U\le U'$ and $nD\le Nd'$, with
\begin{equation}\label{v9:entry:phase-ceiling}
 |H|\le\max_{t=c_-,c_+}
 \frac{F_t+\sqrt{F_t^2+4(Q_t+\eta)\max(0,B_t)}}{2(Q_t+\eta)}.
\end{equation}
\end{lemma}
\begin{proof}
Complete the raw squares first. Since $T'\le T<Nc_-l$, phase pricing
gives $\eta>0$; only now complete the resource squares. The nonnegative
resource contains $n\phi(0)D/6$, proving the weighted bound at actual $n$.
A phase where both resource quadratics are negative is impossible;
replacing $B_t$ by its positive part enlarges the positive root and gives
\eqref{v9:entry:phase-ceiling}. Intersecting valid bounds preserves the
claims; a resource ceiling cannot replace the raw ceiling in its license.
\end{proof}

First license each rho cell's radius by the heat/ray test. Each signed
pass splits beta into 8 closed intervals, each necessary gamma range
into 8, and the common $m$ range into 2. Restore the original beta interval
at every gamma slice before \eqref{v9:entry:mean-contraction}.
A contradictory necessary condition excludes all descendants of that
parent; active plus excluded equals 128 per processed rho cell. Pool upper
bounds by maxima and lower bounds by minima; remove a rho cell only
when every descendant is excluded. Closed seams may duplicate work,
never lose a witness.
For export retain all 100 closed panels $[j/100,(j+1)/100]$; a boundary
law belongs to both incident panels. Use $\theta=1-m^2/(pq)$ with lower $m$
and the full maximum of $pq$ on each cell-panel intersection, including
$1/4$ across $1/2$. Pool all bounds from those same cells.
\begin{condition}[Four signed passes]\label{v9:entry:signed-condition}
Starting from Condition~\ref{v9:entry:two-pass-condition}, all support,
moment, actual-event, Fourier, denominator, signed-curvature and raw-license
gates hold on this complete four-pass cover. Its 11264 nominal boxes and
complete 100-panel export give, on every nonempty panel,
\begin{gather*}
 \mathcal R_c\le2.1995174,\quad nD\le32.20665378,\quad D\le.76682509,\\
 \rho\le71527/48000,\quad|\beta|\le.50549963,\quad|H|\le1.83605680,
 \quad\supp P\subset(-157/50,157/50).
\end{gather*}
Fourteen rho cells survive and 30 probability panels are nonempty.
\end{condition}
\begin{theorem}[Selected coarse entry]\label{v9:entry:theorem}
Under the two finite conditions, every attained pair with $a\ge0$,
$n\ge42$ and true $C_{n-1}\le\CE$ satisfies these bounds and the complete
probability cover, assuming \eqref{v9:entry:structural}.
\end{theorem}
\begin{proof}
Support licenses higher moments, and the actual events and Fourier
appendix bound the full integral. The two coarse passes acquire the gap.
Every signed cell applies only necessary conditions of that same law
and the ordered step. Induction over all four passes preserves the
weighted invariants; closed coverage and outward pooling place every
witness in an incident output panel.
\end{proof}
The numerical conditions remain exact finite obligations. The one-strip
argument removes only downstream panels; this internal cover stays intact.
